\newcommand{\docshorttitle}{SUPPL.\ MEMO P\&A --- \S 425.17(c) COMMERCIAL CARVE-OUT (CGC-25-631801)}
\newcommand{\doctitleblock}{PLAINTIFF'S SUPPLEMENTAL MEMORANDUM OF POINTS AND AUTHORITIES IN FURTHER OPPOSITION TO DEFENDANTS' SPECIAL MOTION TO STRIKE (CODE CIV.\ PROC., \S\S~425.16, 425.17(c))}
\newcommand{\docsubblock}{(Cal.\ Rules of Court, rules 3.1113, 3.1306(c); supplemental to Plaintiff's previously-filed Opposition; cross-noticed for hearings of May 12, 2026 and May 22, 2026, Dept.\ 301)}
% Pleading-paper preamble for APRIL-24-2026 cover sheets and oral-argument
% outlines (CGC-25-631801 and CGC-25-631802). Matches house style from
% APRIL-23-SAC-DEFENSE/court-pdfs/REPLY-SAC-PREAMBLE.tex, simplified for
% single-to-three-page artifacts that do not require TOC/TOA infrastructure.
%
% Cross-hearing caption (CAPTION-801/802.tex): use \CaptionDocTitleBlock,
% \CaptionDocSubBlock, and the crosshearingsschedule environment (defined below)
% so long \doctitleblock text and hearing dates stay inside the 3in minipage.
\documentclass{article}
\usepackage[utf8]{inputenc}
\usepackage{graphicx}
\usepackage{geometry}
    \geometry{top=1in, bottom=2.2in, left=1.0in, right=1.0in, textheight=336pt}
    \setlength{\footskip}{70pt}
\usepackage{ulem}
\usepackage{fancyhdr}
    \renewcommand{\headrulewidth}{0pt}
    \renewcommand{\footrulewidth}{0.4pt}
\usepackage{parskip}
    \setlength{\parskip}{0mm}
\usepackage{quoting}
    \quotingsetup{leftmargin=0.5in, rightmargin=0.5in, vskip=-1.5mm}
    \makeatletter
        \g@addto@macro\quoting\singlespacing
        \g@addto@macro\quoting{\vspace{-2mm}}
    \makeatother
\usepackage{mathptmx}
\usepackage{amssymb}
\renewcommand{\rmdefault}{ptm}
\renewcommand{\normalsize}{\fontsize{12}{12}\selectfont}
\newcommand*{\ptm}{\fontfamily{ptm} \selectfont \fontsize{12}{12}\selectfont}
\usepackage{setspace}
    \doublespacing
\newcommand{\tab}{\hspace*{.0in}}
\newenvironment{tightcenter}{%
    \setlength\topsep{0pt}
    \setlength\parskip{0pt}
    \begin{center}
}{%
    \end{center}
}
\usepackage{tikz}
\usetikzlibrary{calc}
% Pleading-paper vertical double rule at the left margin.
\AddToHook{shipout/background}{%
    \begin{tikzpicture}[remember picture, overlay]
        \draw[line width=0.5pt] ($(current page.north west) + (0.75in,0)$) -- ($(current page.south west) + (0.75in,0)$);
        \draw[line width=0.5pt] ($(current page.north west) + (0.80in,0)$) -- ($(current page.south west) + (0.80in,0)$);
    \end{tikzpicture}%
}
\usepackage[left, pagewise]{lineno}
\setlength{\linenumbersep}{0.35in}
\renewcommand{\linenumberfont}{\normalfont\small}
\usepackage{titlesec}
\usepackage[hidelinks]{hyperref}
\usepackage{bookmark}
\usepackage{textcomp}
\usepackage{longtable}
\usepackage{booktabs}
\usepackage{array}
\usepackage{multirow}
\usepackage{calc}
% --- Cross-hearing caption cover (3in minipage): compact title + 3-column hearing table ---
% Use on first page only; keeps long \doctitleblock / \docsubblock on one sheet with party caption.
\newcommand{\CaptionDocTitleBlock}{%
  \begingroup
  \fontsize{10}{11.5}\selectfont
  \bfseries
  \raggedright
  \setlength{\parskip}{0pt}%
  \noindent\doctitleblock\par
  \endgroup
}
\newcommand{\CaptionDocSubBlock}{%
  \begingroup
  \footnotesize
  \raggedright
  \setlength{\parskip}{0pt}%
  \noindent\docsubblock\par
  \endgroup
}
% #1 = department number only (e.g.\ 301 or 302) for the line ``Dept.\ #1''.
\newenvironment{crosshearingsschedule}[1]{%
  \par\begingroup
  \footnotesize
  \setlength{\tabcolsep}{2.2pt}%
  \renewcommand{\arraystretch}{1.04}%
  \setlength{\parskip}{0pt}%
  \setlength{\parindent}{0pt}%
  \noindent\textbf{Cross-noticed for pending defense hearings (this case, Dept.\ #1):}\par\vspace{2pt}%
  \begin{tabular}{@{}>{\raggedright\arraybackslash}p{0.11\linewidth}>{\raggedright\arraybackslash}p{0.34\linewidth}>{\raggedright\arraybackslash}p{0.53\linewidth}@{}}%
  \textbf{\#} & \textbf{Motion} & \textbf{Date / dept / judge} \\ \hline
}{%
  \end{tabular}\par\endgroup\smallskip
}
% Pandoc pipe-table column widths use \real{#}.
\newcommand{\real}[1]{\pgfmathparse{#1}\pgfmathresult}
\titleformat{\section}{\fontsize{13}{13}\selectfont\normalfont\bfseries\uppercase}{}{0em}{}
\titleformat{\subsection}{\fontsize{12}{12}\selectfont\normalfont\bfseries}{}{0em}{}
\titleformat{\subsubsection}{\normalfont\bfseries}{}{0em}{}
\titlespacing*{\section}{0pt}{6pt}{2pt}
\titlespacing*{\subsection}{0pt}{4pt}{2pt}
\titlespacing*{\subsubsection}{0pt}{2pt}{2pt}
\newcommand{\calrules}[1]{Cal. Rules of Court, rule #1}
\newcommand{\ccp}[1]{Code Civ. Proc., \textsection~#1}
\newcommand{\evidcode}[1]{Evid. Code, \textsection~#1}
\providecommand{\tightlist}{%
  \setlength{\itemsep}{0pt}\setlength{\parskip}{0pt}}
% Footer: page N of total + document short-title.
\usepackage{lastpage}
\pagestyle{fancy}
\fancyhf{}
\fancyfoot[C]{\thepage{} of \pageref{LastPage} \\[1pt] \footnotesize \docshorttitle}
\fancyfoot[R]{}

\begin{document}
% Cross-hearing caption for CGC-25-631801 (Dept. 301, Hon. Van Aken).
\nolinenumbers
\begin{singlespace}
\noindent Franciscus Dylan Rosario \\
7624 Melody Drive \\
Rohnert Park, CA 94928 \\
Tel: 650.488.7510 \\
E: soltrinox@gmail.com \\
Plaintiff, In Pro Per
\end{singlespace}
\vspace*{9mm}
\begin{tightcenter}
\textbf{SUPERIOR COURT OF THE STATE OF CALIFORNIA} \\
\textbf{COUNTY OF SAN FRANCISCO}
\end{tightcenter}
\vspace*{3mm}
\begin{minipage}[t]{3in}
    \raggedright
    \textbf{FRANCISCUS DYLAN ROSARIO,} \\ \textbf{PLAINTIFF}, \\
    \hspace{1cm} v. \\
    \small
    Subhi Abdelhalim; CSAA Insurance Exchange; Phillips, Spallas \& Angstadt LLP; Priya D. Navaratnasingham; Alberto Reyna; Michael R. Chambers; Carbone, Smith \& Koyama; and Does 1 through 50, inclusive;\\
    \normalsize
    \textbf{DEFENDANTS}
    \makebox[3in]{\hrulefill}
\end{minipage}%
\hspace{3mm}%
\begin{minipage}[t]{0.5pt}
    \vspace{0pt}
    \rule{0.5pt}{2.42in}
\end{minipage}%
\hspace{3mm}%
\begin{minipage}[t]{3in}
    \raggedright
    \noindent\textbf{Case No.: CGC-25-631801}\par\smallskip
    \CaptionDocTitleBlock\smallskip
    \CaptionDocSubBlock\smallskip
    \begin{singlespace}
    \setlength{\parindent}{0pt}
    \begin{crosshearingsschedule}{301}
    (1) & MFL / leave (SAC) & \textbf{Apr.\ 30, 2026} \\
    (2) & Anti-SLAPP (\ccp{425.16}) & \textbf{May 12, 2026, 9:30 a.m.};\ Dept.\ 301;\ Hon.\ Christine Van Aken \\
    (3) & Strike (\ccp{436}) & \textbf{May 22, 2026};\ Dept.\ 301;\ Hon.\ Christine Van Aken \\
    (4) & Demurrer & Hrng.\ date \textbf{[TBD]} (reserved) \\
    \end{crosshearingsschedule}
    \noindent Complaint filed: Dec.\ 23, 2025 \\
    \noindent FAC filed: Feb.\ 25, 2026 \\
    \end{singlespace}
\end{minipage}
\vspace*{8mm}
\newpage
\linenumbers


\section*{I. INTRODUCTION AND VEHICLE}

This Supplemental Memorandum of Points and Authorities is respectfully submitted in further opposition to Defendants' special motion to strike under Code of Civil Procedure section 425.16. It is offered as a supplement to --- and not as a substitute for --- Plaintiff's previously-filed Opposition. The vehicle is \calrules{3.1113} (memorandum of points and authorities) read together with \calrules{3.1306(c)} (lodging of supplemental authority for hearings), as the supporting Notice of Lodging Supplemental Authority, the Supplemental Request for Judicial Notice, and the bench briefs filed concurrently in this same cross-hearing package make clear. Plaintiff incorporates the previously-filed Opposition by reference and submits this supplemental memorandum to focus the Court's attention on the structurally prior question raised by Code of Civil Procedure section 425.17(c): whether the special motion to strike procedure is available to Defendants at all.

As shown below, it is not. The moving Defendants are either the commercial insurer itself (CSAA Insurance Exchange) or panel counsel and individual lawyers retained by, or acting in alignment with, the insurer in furtherance of CSAA's commercial claim-handling and litigation posture. Plaintiff's causes of action arise from representations of fact about the insurer's claim handling and from conduct made in the course of delivering the insurance product to Plaintiff and to decision-makers in a position to influence Plaintiff's customer-side claim. Both the commercial-actor element and the conduct-and-audience elements of \ccp{425.17(c)} are satisfied. Neither Civil Code section 47(b) nor Business and Professions Code section 6128(a) supplies a route around the carve-out. The special motion to strike must be denied at the gateway without reaching the merits-stage \ccp{425.16(b)(1)} ``probability of prevailing'' inquiry.

The Supplemental Request for Judicial Notice filed concurrently reproduces the verbatim text of \ccp{425.17}, \ccp{425.16(e), (g), (i)}, \ccp{904.1(a)(13)}, Civil Code \S~47(b), Business and Professions Code \S~6128(a), Insurance Code \S~790.03(h), and 10 California Code of Regulations \S~2695.1. A companion Supplemental Memorandum filed concurrently in this case (the ``Customer-Doctrine Supplemental'') addresses the discrete question whether a third-party claimant is a ``customer'' within the meaning of section 425.17(c)(2) and whether panel counsel are themselves operating within the regulated commercial channel; that question is treated only summarily here.

\section*{II. STATEMENT OF FACTS GERMANE TO THE \S~425.17(c) GATEWAY}

The moving Defendants in CGC-25-631801 include CSAA Insurance Exchange (a California-licensed insurer) and several panel counsel and individual lawyers retained by, or acting in alignment with, CSAA in furtherance of the insurer's commercial litigation posture. Plaintiff's pleaded causes of action arise from: (a) representations of fact about CSAA's claims handling and the status and conclusions of its investigation; (b) communications regarding the basis for denial; and (c) representations concerning the existence or non-existence of factual predicates that were communicated to Plaintiff and to the court before trial. Each set of communications was made in the course of delivering the insurance product (as to CSAA) or in the course of delivering the insurance defense product retained and paid for by CSAA (as to panel counsel). The intended audience at every stage was an actual or potential decision-maker influencing the outcome of the customer's claim --- first Plaintiff as the insured, then the regulator, and ultimately the court.

\section*{III. LEGAL STANDARD --- \S~425.17(c) IS A STRUCTURAL GATEWAY}

\subsection*{A. The Carve-Out Is Categorical and Sequentially Prior to \S~425.16.}

The Legislature enacted Code of Civil Procedure section 425.17 in 2003 because it determined that \ccp{425.16} was being misused. \ccp{425.17(a)} declares: ``[T]here has been a disturbing abuse of Section 425.16, the California Anti-SLAPP Law, which has undermined the exercise of the constitutional rights of freedom of speech and petition for the redress of grievances, contrary to the purpose and intent of Section 425.16.'' Subdivision (c) is the Legislature's response to one identifiable category of abuse: special motions to strike filed by commercial actors --- including, by name, those engaged in the business of insurance --- to defeat causes of action that arise from their own commercial representations or from conduct undertaken in the course of delivering their goods or services to actual or potential customers.

The California Supreme Court has confirmed in \textit{Simpson Strong-Tie Co.\ v.\ Gore} (2010) 49 Cal.4th 12, and the Court of Appeal has reaffirmed in \textit{JAMS, Inc.\ v.\ Superior Court} (2016) 1 Cal.App.5th 984 and \textit{Stewart v.\ Rolling Stone LLC} (2010) 181 Cal.App.4th 664, that the carve-out is categorical: when its elements are met, \ccp{425.16} ``does not apply'' at all. The Court does not weigh competing interests; it declines to extend the special motion to strike procedure to the category of disputes the Legislature has removed from that mechanism. \textit{Club Members for an Honest Election v.\ Sierra Club} (2008) 45 Cal.4th 309 supplies the sequencing rule: where \ccp{425.17} is invoked, the trial court must address the carve-out before any \ccp{425.16} analysis. This sequencing is not optional.

\subsection*{B. The Two-Element Test of \S~425.17(c).}

The two elements of \ccp{425.17(c)} are straightforward. \textbf{First}, the defendant must be ``primarily engaged in the business of selling or leasing goods or services,'' with insurance specifically enumerated. \textbf{Second}, the cause of action must arise from a statement or conduct that consists either of representations of fact about the defendant's business operations, goods, or services, made for the purpose of obtaining approval for, promoting, or securing a commercial transaction, \textbf{or} of conduct made in the course of delivering those goods or services. In addition, \ccp{425.17(c)(2)} requires that the intended audience be (i) an actual or potential buyer or customer; (ii) a person likely to repeat the statement to or otherwise influence such a customer; or (iii) alternatively, that the conduct have arisen out of or within the context of a regulatory approval process, proceeding, or investigation. Each audience branch is independently sufficient.

\section*{IV. ARGUMENT --- THE \S~425.17(c) CARVE-OUT REQUIRES DENIAL OF THE SPECIAL MOTION TO STRIKE}

\subsection*{A. CSAA Is the Paradigmatic \S~425.17(c) Defendant; Panel Counsel Operate Within the Same Commercial Channel.}

CSAA is a California-licensed insurer. Insurance is one of the named industries in the text of \ccp{425.17(c)}. The first element is satisfied as to CSAA as a matter of law.

\ccp{425.17(c)} is not limited to the commercial actor in name; its focus is on the \textbf{cause of action} and what it \textbf{arises from}, not on the formal corporate identity of the speaker. Where panel counsel and counsel's firm are retained by an insurer to deliver --- as part of the insurance product itself --- defense services that the policy promises to the insured, those attorneys are themselves operating within the commercial channel that \ccp{425.17(c)} regulates. Where Plaintiff's claims against panel counsel arise from representations of fact about CSAA's claim handling, about the status of an investigation, about the basis for a denial, or about the existence or non-existence of factual predicates that were communicated to the court before trial, those representations were made in the course of delivering the insurance defense product to the insurer-client and, derivatively, were made about the insurer's operations within the meaning of \ccp{425.17(c)(1)}.

\subsection*{B. The Intended-Audience Element Is Independently Satisfied Under Both Available Branches.}

The intended audience for the operative communications was, at every stage, an actual or potential decision-maker influencing the outcome of the customer's claim --- first Plaintiff as the insured, then the regulator, and ultimately the court. That falls within both branches of \ccp{425.17(c)(2)}: (i) the ``person likely to repeat the statement to, or otherwise influence'' the customer branch; and (ii) the ``regulatory approval process, proceeding, or investigation'' branch, given that Insurance Code \S~790.03(h) and 10 CCR \S~2695.1 et seq.\ directly govern misrepresentations to claimants, the duty to adopt reasonable investigation standards, and the obligation to attempt prompt, fair, and equitable settlement.

\subsection*{C. Civil Code \S~47(b) and the ``Petitioning'' Gloss of \S~425.16(e) Do Not Override the Carve-Out.}

Defendants rely on Civil Code \S~47(b) and on the line of cases extending the litigation privilege to communications ``in furtherance of'' judicial proceedings. As \textit{Action Apartment Ass'n v.\ City of Santa Monica} (2007) 41 Cal.4th 1232 confirms, \S~47(b) is, in many respects, coextensive with \ccp{425.16(e)} for certain deceit-based predicates. That coextensiveness cuts both ways: if the conduct does not satisfy the petitioning concept underlying \ccp{425.16(e)}, it likewise cannot be sheltered by \S~47(b) for purposes of evading \ccp{425.17(c)}. And \S~47(b) itself contains carve-outs --- including subdivision (b)(2) for communications made in furtherance of intentional destruction or alteration of physical evidence, and subdivision (b)(3) for knowing concealment of insurance policies in a judicial proceeding.

Business and Professions Code \S~6128(a) supplies an additional, independent prohibition: every attorney is guilty of a misdemeanor who is ``guilty of any deceit or collusion, or consents to any deceit or collusion, with intent to deceive the court or any party.'' Section 6128(a) is not a cause of action that displaces other privileges; it is, however, an authoritative statement that California law does not regard attorney deceit aimed at the court as protected speech. Read alongside \ccp{425.17(c)}, the result is unambiguous: an insurer's panel counsel cannot, by labeling pre-trial commercial representations as ``advocacy,'' convert deceit aimed at the tribunal into anti-SLAPP-protected petitioning activity.

\subsection*{D. \textit{Flatley} Independently Closes the \S~425.16 Gateway as to Conduct Conclusively Shown to Be Unlawful.}

\textit{Flatley v.\ Mauro} (2006) 39 Cal.4th 299, reaffirmed in \textit{Lefebvre v.\ Lefebvre} (2011) 199 Cal.App.4th 696, holds that conduct conceded or shown as a matter of law to be illegal is not protected by \ccp{425.16} in the first instance. To the extent Defendants' special motion depends on conduct that, on an uncontroverted record, violates Insurance Code \S~790.03(h), Business and Professions Code \S~6128(a), or constitutes extrinsic fraud, the \textit{Flatley} gateway closes regardless of how the conduct is labeled.

\subsection*{E. Extrinsic Fraud Is Not ``Protected Speech.''}

Extrinsic fraud is, by definition, fraud that prevents a party from fully and fairly presenting his case --- including misrepresentations about whether an investigation has occurred, false assurances that govern a litigant's pre-trial conduct, and concealment of facts the opposing party was entitled to know. The litigation privilege has never been understood to immunize extrinsic fraud; if it did, the doctrine of extrinsic fraud would be a dead letter. Where panel counsel's challenged statements are alleged to have been made and acted upon \textbf{prior to trial}, with the consequence of impairing Plaintiff's ability to present his case, those statements are not ``protected speech'' under any reading of \ccp{425.16(e)}. They are, instead, the very conduct \ccp{425.17(c)} was enacted to keep within the ordinary civil docket. \textit{Rusheen v.\ Cohen} (2006) 37 Cal.4th 1048 draws a clear line between communicative acts integral to a judicial proceeding and the underlying non-communicative or pre-litigation conduct that gives rise to a cause of action; the gravamen of Plaintiff's claims is the latter.

\subsection*{F. The ``Concealment of the Failure to Investigate'' Theory Is Squarely Within \S~425.17(c)(1).}

Plaintiff's pleaded theory --- that the insurer communicated a denial premised on a represented conclusion of investigation when in fact no adequate investigation had occurred, and that panel counsel carried that representation forward --- is structurally a \S~790.03(h)(1) and (h)(3) theory. The Court need not decide on this motion whether each statutory element has been violated; it need only recognize that conduct alleged to have violated \S~790.03(h) is, by definition, conduct ``in the course of delivering'' insurance services and therefore squarely within \ccp{425.17(c)(1)}.

\subsection*{G. The Cross-Noticed \S~436 Motion Does Not Rescue the Defense.}

Defendants have cross-noticed a motion under \ccp{436} that incorporates by reference the ``everything is petitioning'' rhetoric of their anti-SLAPP papers. \textit{PH II, Inc.\ v.\ Superior Court} (1995) 33 Cal.App.4th 1680, 1682--83, holds that a \ccp{436} motion is not a backdoor demurrer; by parity of reasoning, it is not a backdoor anti-SLAPP motion either. The underlying conduct remains commercial-insurer claims handling and counsel's role in delivering that product, not protected petitioning.

\subsection*{H. \S~425.17(e) --- No Interlocutory Appeal, No Stay, No Writ.}

If this Court denies the special motion to strike on \ccp{425.17} grounds --- as it must on this record --- the denial is not subject to interlocutory appeal under \ccp{904.1(a)(13)}. \ccp{425.17(e)} is explicit. The Legislature concluded that the ordinary anti-SLAPP appeal route is itself a vehicle for abuse where the underlying claim was never within \ccp{425.16} to begin with.

\section*{V. CONCLUSION}

In CGC-25-631801, the moving Defendants are either the commercial insurer itself or counsel acting in the course of delivering the insurance product. The conduct from which Plaintiff's causes of action arise consists of representations of fact about the insurer's operations and claims, made to Plaintiff as customer or to decision-makers in a position to influence his customer-side claim, and made within the regulatory framework of Insurance Code \S~790.03(h) and the Fair Claims Settlement Practices Regulations. Defendants' ``freedom of speech'' framing depends on a re-characterization that \ccp{425.17(c)} was enacted to forbid, and Civil Code \S~47(b) does no independent work because the litigation privilege does not immunize extrinsic fraud, does not displace \ccp{425.17(c)}'s structural sequencing, and does not override Business and Professions Code \S~6128(a)'s prohibition on attorney deceit aimed at the court or any party. Plaintiff respectfully requests that the Court deny the special motion to strike at the gateway under \ccp{425.17(c)}, without reaching the second prong of the anti-SLAPP analysis, and proceed with the ordinary civil docket in CGC-25-631801.

\input{SIGBLOCK.tex}
\end{document}
